\documentclass[11pt,letterpaper]{article}

\usepackage[T1]{fontenc}
\usepackage[utf8]{inputenc}
\usepackage[spanish,es-noshorthands,es-tabla]{babel}
\usepackage{lmodern}
\usepackage{microtype}
\usepackage[margin=2.3cm]{geometry}
\usepackage{amsmath,amssymb}
\usepackage{booktabs}
\usepackage{graphicx}
\usepackage{enumitem}
\usepackage{textcomp}
\usepackage{xcolor}
\usepackage{hyperref}
\usepackage[capitalise,nameinlink,spanish]{cleveref}
\usepackage{tabularx}
\usepackage[section]{placeins}
\crefname{table}{Tabla}{Tablas}
\Crefname{table}{Tabla}{Tablas}
\crefname{section}{Sección}{Secciones}
\Crefname{section}{Sección}{Secciones}
\crefname{subsection}{Sección}{Secciones}
\Crefname{subsection}{Sección}{Secciones}
\hypersetup{colorlinks=true,linkcolor=blue!50!black,urlcolor=blue!60!black,citecolor=blue!50!black}
\setlist{itemsep=2pt,topsep=4pt}

% generado por analisis/graficos.py; no editar a mano
\newcommand{\resCasos}{1\,140}
\newcommand{\resIguales}{1\,110}
\newcommand{\resIgualesEsc}{280/280}
\newcommand{\resIgualesPor}{600/600}
\newcommand{\resIgualesCol}{150/180}
\newcommand{\resIgualesVar}{80/80}
\newcommand{\resUniDInc}{152}
\newcommand{\resUniDRan}{172}
\newcommand{\resUniDCociente}{1.14}
\newcommand{\resUniIInc}{134}
\newcommand{\resUniIRan}{155}
\newcommand{\resUniICociente}{1.16}
\newcommand{\resCirDInc}{157}
\newcommand{\resCirDRan}{1\,479}
\newcommand{\resCirDCociente}{9.39}
\newcommand{\resCirIInc}{128}
\newcommand{\resCirIRan}{1\,051}
\newcommand{\resCirICociente}{8.22}
\newcommand{\resIntDoubleUniInc}{0.88}
\newcommand{\resIntDoubleCirInc}{0.81}
\newcommand{\resIntDoubleUniRan}{0.90}
\newcommand{\resIntDoubleCirRan}{0.71}
\newcommand{\resIncNormMin}{7.3}
\newcommand{\resIncNormMax}{7.7}
\newcommand{\resCirRanNormMin}{32}
\newcommand{\resCirRanNormMax}{74}
\newcommand{\resCostoPequeMin}{24}
\newcommand{\resCostoPequeMax}{36}
\newcommand{\resCostoGrandeMin}{31}
\newcommand{\resCostoGrandeMax}{70}
\newcommand{\resRporNMin}{4.1}
\newcommand{\resRporNMax}{4.3}
\newcommand{\resRuidoMin}{3}
\newcommand{\resRuidoMax}{29}
\newcommand{\resCorrUniD}{0.70}
\newcommand{\resCorrCirD}{0.00}
\newcommand{\resCorrUniI}{0.74}
\newcommand{\resCorrCirI}{-0.08}
\newcommand{\resUniCocMinD}{0.90}
\newcommand{\resUniCocMinND}{50\,000}
\newcommand{\resUniVentanaD}{20\,000--200\,000}
\newcommand{\resUniCocMinI}{0.84}
\newcommand{\resUniCocMinNI}{100\,000}
\newcommand{\resUniVentanaI}{20\,000--500\,000}
\newcommand{\resUniH}{36}
\newcommand{\resUniHEsperado}{37}
\newcommand{\resReasigTodos}{22\,840\,921}
\newcommand{\resReasigTodosCota}{28\,785\,453}
\newcommand{\resReasigTodosPorN}{22.8}
\newcommand{\resIntTodosH}{102\,550}
\newcommand{\resIntTodosHcienmil}{57\,840}
\newcommand{\tablaCruce}{%
double & 10\,000 & 1.09 & 3.98 & no se cruzan \\
double & 100\,000 & 0.93 & 4.72 & 2.8\,\% \\
double & 1\,000\,000 & 1.21 & 9.92 & no se cruzan \\
int & 10\,000 & 1.12 & 4.00 & no se cruzan \\
int & 100\,000 & 1.02 & 5.24 & 2.2\,\% \\
int & 1\,000\,000 & 1.03 & 7.35 & no se cruzan \\
}
\newcommand{\tablaDescomposicion}{%
double & 10\,000 & 13.3 & 4.31 & 0.32 & 7.8 & 26.3 & 3.4 & 1.09 \\
double & 100\,000 & 16.6 & 4.20 & 0.25 & 7.3 & 27.0 & 3.7 & 0.93 \\
double & 1\,000\,000 & 19.9 & 4.08 & 0.20 & 7.4 & 43.6 & 5.9 & 1.21 \\
int & 10\,000 & 13.3 & 4.31 & 0.32 & 7.2 & 24.8 & 3.5 & 1.12 \\
int & 100\,000 & 16.6 & 4.20 & 0.25 & 6.6 & 26.6 & 4.0 & 1.02 \\
int & 1\,000\,000 & 19.9 & 4.08 & 0.20 & 6.5 & 33.0 & 5.1 & 1.03 \\
}
\newcommand{\tablaVarianza}{%
double & uniforme & Incremental & 155.1 & 12.5 & 8.1\,\% & -- & -- \\
double & uniforme & Random Incremental & 199.8 & 34.8 & 17.4\,\% & 4\,227\,082 & 19.99\,\% \\
double & 100\,\% en la cerradura & Incremental & 169.3 & 12.4 & 7.3\,\% & -- & -- \\
double & 100\,\% en la cerradura & Random Incremental & 1804.3 & 202.3 & 11.2\,\% & 22\,918\,649 & 2.23\,\% \\
int & uniforme & Incremental & 134.4 & 8.2 & 6.1\,\% & -- & -- \\
int & uniforme & Random Incremental & 146.6 & 24.8 & 16.9\,\% & 4\,227\,082 & 19.99\,\% \\
int & 100\,\% en la cerradura & Incremental & 123.0 & 8.5 & 6.9\,\% & -- & -- \\
int & 100\,\% en la cerradura & Random Incremental & 990.6 & 97.2 & 9.8\,\% & 20\,994\,650 & 2.43\,\% \\
}
\newcommand{\tablaColinealidad}{%
cuadrado, 0\textdegree & int & 10\,000 & 100\,\% & 4 & 4 \\
cuadrado, 0\textdegree & int & 100\,000 & 100\,\% & 4 & 4 \\
cuadrado, 0\textdegree & int & 1\,000\,000 & 100\,\% & 4 & 4 \\
cuadrado, 0\textdegree & double & 10\,000 & 100\,\% & 4 & 4 \\
cuadrado, 0\textdegree & double & 100\,000 & 100\,\% & 4 & 4 \\
cuadrado, 0\textdegree & double & 1\,000\,000 & 100\,\% & 4 & 4 \\
\midrule
cuadrado, 30\textdegree & int & 10\,000 & 100\,\% & 25--35 & 25--35 \\
cuadrado, 30\textdegree & int & 100\,000 & 100\,\% & 36--47 & 36--47 \\
cuadrado, 30\textdegree & int & 1\,000\,000 & 100\,\% & 41--58 & 41--58 \\
cuadrado, 30\textdegree & double & 10\,000 & 0\,\% & 21--29 & 10--16 \\
cuadrado, 30\textdegree & double & 100\,000 & 0\,\% & 39--55 & 14--21 \\
cuadrado, 30\textdegree & double & 1\,000\,000 & 0\,\% & 46--67 & 19--26 \\
\midrule
circunferencia & int & 10\,000 & 100\,\% & 1\,000 & 1\,000 \\
circunferencia & int & 100\,000 & 100\,\% & 10\,000 & 10\,000 \\
circunferencia & int & 1\,000\,000 & 100\,\% & 57\,752--58\,092 & 57\,752--58\,092 \\
circunferencia & double & 10\,000 & 100\,\% & 1\,000 & 1\,000 \\
circunferencia & double & 100\,000 & 100\,\% & 10\,000 & 10\,000 \\
circunferencia & double & 1\,000\,000 & 100\,\% & 100\,000 & 100\,000 \\
}

% generado por analisis/exactitud.py; no editar a mano
\newcommand{\tablaExactitud}{%
10\,000 & 17--23 & 0/10 & 7.0 & 11.0 & 0/10 & 9.1 & 1.3 \\
100\,000 & 26--38 & 0/10 & 12.5 & 27.0 & 0/10 & 15.1 & 1.9 \\
}


\newcommand{\inc}{\textsf{Incremental}}
\newcommand{\ran}{\textsf{Random Incremental}}
\newcommand{\code}[1]{\texttt{#1}}

\title{Tarea 2: Cerradura convexa\\[2pt]
  \large CC5502 Geometría Computacional --- ANALISIS}
\author{Matias Saavedra}
\date{Octubre 2026}

\begin{document}
\maketitle

\section{Objetivo y preguntas}

Se implementaron dos algoritmos de cerradura convexa en C++, sobre las clases
\code{Punto}, \code{Vector} y \code{Poligono} de la Tarea~1: el \inc{} (ordenar,
construir el \emph{upper hull} y devolverse construyendo el \emph{lower hull}) y
el \ran{} de la sección~2 de las notas de CMU que cita el
enunciado~\cite{cmu}. Los dos son $O(n\log n)$ (el segundo en esperanza), así
que el análisis asintótico no dice cuál conviene: lo deciden las constantes y
cómo depende cada uno de la forma de la entrada. Siguiendo la Lectura~2
\cite{moret}, las preguntas se fijaron antes de medir:

\begin{enumerate}[label=\textbf{P\arabic*}]
  \item ¿Crece el tiempo de cada algoritmo como $n\log n$?
  \item Con puntos aleatorios uniformes, ¿cuál es más rápido y por cuánto?
  \item ¿Cómo depende el tiempo del porcentaje de puntos que forman la
        cerradura? ¿Hay un porcentaje desde el cual conviene uno u otro, y se
        puede determinar?
  \item ¿Qué cambia entre coordenadas enteras (\code{int}) y reales
        (\code{double})?
  \item ¿Coinciden las cerraduras cuando hay puntos colineales en el borde, y
        cuando no los hay?
\end{enumerate}


\section{Implementación}

\subsection{Estructura}

El código está en los headers de \code{include/}, porque todo es template
(\code{T} = \code{int} o \code{double}). Siguiendo el patrón Strategy,
\code{EstrategiaCerradura<T>} declara el método
\begin{center}
\code{virtual Poligono<T> apply(const std::vector<Punto<T>> \&nube) = 0;}
\end{center}
y de ella heredan \code{IncrementalAlgorithm} y \code{RandomIncrementalAlgorithm}. Los generadores de nubes siguen el mismo
patrón: \code{GeneradorPuntos<T>} declara \code{generar(n, rng)} y de él
heredan \code{GeneradorUniforme}, \code{GeneradorCircunferencia} y
\code{GeneradorCuadrado}. El programa del experimento solo recorre listas de
estrategias y generadores.

Los dos algoritmos usan un único predicado geométrico, \code{orientacion(a, b,
c)}: el signo del producto cruz $(b-a)\times(c-a)$, calculado con
\code{Vector::productoCruz} de la Tarea~1. Con \code{T} entero las
coordenadas se promueven a \code{long long} antes de restar, y el signo es
exacto mientras $|x|,|y|<2^{30}$; con \code{double} el producto se redondea.
De la Tarea~1 se agregaron \code{Punto::operator<} (orden lexicográfico), un
constructor vacío de \code{Poligono}, \code{Poligono::normalizar} y
\code{Poligono::operator==}.

\subsection{Incremental}

Ordena los puntos por $x$ (y por $y$ a igual $x$), elimina los repetidos y
recorre la lista de izquierda a derecha construyendo el \emph{upper hull}: como
el punto nuevo es el extremo derecho de los ya vistos, siempre es vértice de su
cerradura, y basta agregarlo al final y sacar los últimos mientras los tres
últimos no formen una vuelta a la derecha. Luego se devuelve de derecha a
izquierda con la misma regla para el \emph{lower hull}
(\cite{deberg}, sección~1.1). Cada punto entra y sale de la cadena a lo más una
vez por recorrido, así que tras el orden $O(n\log n)$ el resto es $O(n)$, y el
costo \emph{no depende} de cuántos puntos quedan en la cerradura.

\subsection{Random Incremental}

Baraja los puntos (\code{std::shuffle} con \code{std::mt19937\_64} y una
semilla fija), toma un triángulo no degenerado y agrega los demás puntos en ese
orden. La dificultad es saber, para cada punto nuevo, si está dentro y qué
aristas ve. Como en \cite{cmu}, se fija un punto $C$ interior (el centroide del
triángulo inicial; la cerradura solo crece, así que $C$ sigue dentro) y cada
punto que falta por insertar guarda la arista que cruza el rayo $C\to q$. Esa
arista es visible desde $q$ exactamente cuando $q$ está fuera, de modo que
\emph{$q$ está fuera si y solo si está estrictamente a la derecha de su arista}:
una sola orientación. Cada arista guarda, a su vez, la ``cubeta'' de puntos que
la apuntan. Al insertar $q$:

\begin{enumerate}
  \item se recorre la cerradura desde la arista de $q$ hacia ambos lados
        mientras las aristas sean visibles, hasta los vértices $l$ y $r$;
  \item se sacan las cubetas de las aristas visibles y la cadena $l\ldots r$ se
        reemplaza por $l\to q\to r$;
  \item cada punto de esas cubetas cruza ahora $l\to q$ o $q\to r$, según de qué
        lado del rayo $C\to q$ esté, y pasa a esa cubeta; si ya no ve su nueva
        arista, está dentro de la cerradura y se descarta para siempre.
\end{enumerate}

Por análisis hacia atrás, el número esperado de reasignaciones es a lo más
$2nH_n\approx 2n\ln n$ para cualquier entrada \cite{cmu}: la esperanza es sobre
el barajado, no sobre los datos. El descarte del paso~3 no está en las notas;
no cambia la cota pero hace que el trabajo dependa de cuántos puntos siguen
fuera de la cerradura, que es lo que mide la pregunta~P3. El método
\code{getReasignaciones()} entrega el contador para el experimento.

\paragraph{Calidad de la implementación.} Moret advierte que una implementación
descuidada cambia las conclusiones \cite{moret}. Se midieron tres versiones del
\ran{} con la misma nube de $n=10^6$ puntos
(\code{resultados/versiones\_random\_incremental.txt}). Guardar las cubetas
como listas enlazadas a través de un arreglo por punto tardaba 1\,997\,ms con
todos los puntos en la cerradura; agrupar los campos de cada punto en un
\code{struct} casi no cambió el tiempo (1\,923\,ms), y guardar cada cubeta en un
\code{std::vector<int>} contiguo lo bajó a 1\,624\,ms (con $p=10$\,\%: 371, 445
y 334\,ms). Se usó esta última, que además es la más simple. El perfil
(\code{valgrind --tool=callgrind}) muestra que el costo restante son las
operaciones mismas, unas 125 instrucciones por reasignación (dos orientaciones,
la cubeta nueva y su crecimiento), más la latencia de memoria cuando $n$ ya no
cabe en el caché.

\subsection{Supuestos y decisiones}
\label{sec:supuestos}

\begin{itemize}
  \item \textbf{Salida canónica.} Ambos algoritmos entregan solo los vértices
        extremos (un punto colineal sobre una arista no es vértice, y los
        repetidos aparecen una vez), en sentido antihorario y desde el menor
        vértice en orden lexicográfico. Así ``cerraduras idénticas'' es
        comparar vértice a vértice (\code{Poligono::operator==}). Con todos los
        puntos colineales la cerradura es el segmento entre los extremos.
  \item \textbf{Colineales en el \inc.} Se saca el último punto de la cadena
        cuando la orientación es $\ge 0$, no solo $>0$.
  \item \textbf{Colineales en el \ran.} La prueba ``$q$ está fuera'' es estricta
        (un punto sobre el borde está dentro), pero el recorrido del paso~1 no:
        también elimina una arista cuya recta pasa por $q$. Si $q$ está en la
        prolongación de una arista, el vértice que comparten deja de ser
        extremo; con un recorrido estricto quedaría como vértice colineal y las
        cerraduras no coincidirían. Las notas suponen que no hay tres puntos
        colineales; esta regla es lo que hace falta sin ese supuesto.
  \item \textbf{Triángulo inicial.} Se toma el primer punto distinto del primero
        y luego el primero no colineal con ambos: con coordenadas enteras o
        colineales los tres primeros puntos del barajado pueden no servir.
  \item \textbf{$C$ exacto.} El centroide $C=(a+b+c)/3$ no tiene coordenadas
        enteras, y un triángulo entero puede no tener ningún punto entero en su
        interior. Las orientaciones respecto de $C$ se calculan con $3p-(a+b+c)$,
        que es $3$ veces el vector $C\to p$. Con \code{T} entero esto es exacto
        en \code{long long} si $|x|,|y|<2^{28}$; los generadores usan
        $|x|,|y|\le 1.42\cdot 10^{8}<2^{28}$.
  \item \textbf{Aritmética inexacta.} Con \code{double} no se garantiza que
        ambos algoritmos coincidan. El recorrido del paso~1 tiene topes para no
        dar la vuelta completa a la cerradura si los signos se contradicen.
\end{itemize}

\subsection{Generadores}

Los tres trabajan en \code{double}, con coordenadas de hasta $R=10^8$, y con
\code{T} entero redondean al entero más cercano: con la misma semilla, la nube
\code{int} es la nube \code{double} redondeada, y la única diferencia entre los
dos tipos es el redondeo.

\begin{itemize}
  \item \code{GeneradorUniforme}: $n$ puntos uniformes en $[-R,R]^2$. Su
        cerradura tiene en esperanza unos $\tfrac{8}{3}\ln n$ vértices
        \cite{renyi}, unos \resUniHEsperado{} para $n=10^6$ (se midió una
        mediana de \resUniH): es el caso de un porcentaje muy pequeño en la
        cerradura.
  \item \code{GeneradorCircunferencia(p)}: $k=\max(3,pn)$ puntos equiespaciados
        sobre la circunferencia de radio $R$, con un desfase aleatorio, y el
        resto uniforme en el disco inscrito en ese $k$-ágono, de radio
        $0.999\,R\cos(\pi/k)$. Los interiores quedan estrictamente dentro del
        polígono, y no solo del círculo: si se sortearan en el disco de radio
        $R$, los que caen entre una arista y el arco serían vértices y el
        porcentaje real sería mayor que el pedido. Así, con \code{double} la
        cerradura tiene exactamente $k$ vértices, sin colineales.
  \item \code{GeneradorCuadrado(p, $\theta$)}: las 4 esquinas y $\max(0,pn-4)$
        puntos más sobre los lados del cuadrado $[-R,R]^2$ rotado en $\theta$, y
        el resto dentro. Con $\theta=0$ los lados son horizontales y verticales
        y la colinealidad es exacta en ambos tipos; con $\theta=30^\circ$ los
        puntos están sobre la recta solo en los reales: en \code{double} quedan a
        un redondeo de ella, y en \code{int} a medio entero.
\end{itemize}

\section{Diseño experimental}

\begin{table}[!htbp]
\centering
\caption{Partes del experimento. Cada caso se mide con \code{int} y con
  \code{double}.}
\label{tab:diseno}
\small
\begin{tabularx}{\linewidth}{@{}lXll@{}}
\toprule
Parte & Generador & $n$ & Repeticiones \\
\midrule
Escalamiento & uniforme; circunferencia 100\,\%
  & $1,2,5\cdot 10^{4}$; $1,2,5\cdot 10^{5}$; $10^{6}$ & 10 nubes \\
Porcentaje & circunferencia con $p$ = 0.1, 0.2, 0.5, 1, 2, 5, 10, 20, 50 y 100\,\%
  & $10^{4}$, $10^{5}$, $10^{6}$ & 10 nubes \\
Colinealidad & cuadrado $0^\circ$ y $30^\circ$; circunferencia ($p=10$\,\%)
  & $10^{4}$, $10^{5}$, $10^{6}$ & 10 nubes \\
Varianza & uniforme; circunferencia 100\,\% & $10^{6}$ & 1 nube, 20 semillas \\
\bottomrule
\end{tabularx}
\end{table}

En total son \resCasos{} casos (\cref{tab:diseno}); en cada uno se genera una
nube y se miden ambos algoritmos sobre esa misma nube.

\begin{itemize}
  \item \textbf{Medidas.} El tiempo de \code{apply} (que incluye la copia de la
        entrada y el orden o el barajado; no incluye generar la nube), el
        tamaño $h$ de cada cerradura, si las dos son idénticas y, para el
        \ran, el número de reasignaciones. Esta última es una medida
        estructural \cite{moret}: se reproduce exactamente y explica el tiempo.
  \item \textbf{Tiempo.} Con \code{std::chrono::steady\_clock}. Cada medición
        repite \code{apply} hasta juntar al menos 50\,ms y divide, para que la
        resolución del reloj no domine con $n$ pequeño. Se reporta la mediana
        de las 10 nubes y, en los gráficos, la banda entre los cuartiles.
  \item \textbf{Confusores controlados.} Ambos algoritmos sobre la misma nube; el
        orden de medición se alterna entre nubes; el \ran{} libera su memoria
        de trabajo al terminar, como el \inc, para que las repeticiones
        reserven memoria igual que la primera; el proceso se fija a un núcleo
        (\code{taskset -c 3}) en una máquina sin otra carga.
  \item \textbf{Variabilidad.} La parte ``varianza'' mide una misma nube con 20
        semillas del barajado: separa la variabilidad propia del \ran{} de la
        que viene de cambiar de nube.
  \item \textbf{Reproducibilidad.} Las semillas son fijas
        (\code{std::mt19937\_64}; la nube $i$ de tamaño $n$ usa la semilla
        $1\,000\,003\,i+n$ y el algoritmo la semilla $i+1$). $h$, las
        reasignaciones y las coincidencias se reproducen exactamente.
  \item \textbf{Máquina.} Intel Core i7-9750H (6 núcleos, 2.6\,GHz base y hasta
        4.5\,GHz, caché L3 de 12\,MiB), 16\,GB de RAM, Linux 7.2.9 (Arch), g++
        16.2.1 con \code{-O3 -DNDEBUG} (\code{Release} de CMake), C++17. El
        regulador de frecuencia (\code{powersave}, con turbo) no se fijó: es la
        principal fuente de ruido en los tiempos.
\end{itemize}


\section{Resultados}

\subsection{Las cerraduras coinciden}

En \resIguales{} de los \resCasos{} casos las dos cerraduras son idénticas:
\resIgualesEsc{} en el escalamiento, \resIgualesPor{} en el porcentaje,
\resIgualesVar{} en la varianza y \resIgualesCol{} en la colinealidad. Las 30
diferencias son todas del mismo caso, el cuadrado rotado en $30^\circ$ con
\code{double}, y se analizan en la \cref{sec:colineal}.

\subsection{Escalamiento con \texorpdfstring{$n$}{n} (P1)}

\begin{figure}[!htbp]
\centering
\includegraphics[width=\linewidth]{figuras/escalamiento.pdf}\\[4pt]
\includegraphics[width=\linewidth]{figuras/escalamiento_normalizado.pdf}
\caption{Arriba: tiempo de \code{apply} contra $n$ (mediana de 10 nubes y banda
  intercuartil). Abajo: el mismo tiempo dividido por $n\log_2 n$; una curva
  horizontal es un crecimiento $n\log n$.}
\label{fig:escalamiento}
\end{figure}

La \cref{fig:escalamiento} confirma el $n\log n$ del \inc: entre $n=10^4$ y
$10^6$ tarda entre \resIncNormMin{} y \resIncNormMax{}\,ns por $n\log_2 n$ con
puntos uniformes, y lo mismo con todos los puntos en la cerradura: su costo no
depende de $h$.

El \ran{} con todos los puntos en la cerradura es el peor caso de su análisis.
La cantidad de operaciones sí crece como $n\log n$: con $n=10^6$ hace
\resReasigTodos{} reasignaciones ($\resReasigTodosPorN\,n$), bajo la cota
$2nH_n=\resReasigTodosCota$, y en la \cref{fig:reasignaciones} las
reasignaciones por punto siguen a $2\ln n$. Pero el tiempo por $n\log_2 n$
sube de \resCirRanNormMin{} a \resCirRanNormMax{}\,ns: lo que crece es el
costo de \emph{cada} reasignación, de \resCostoPequeMin--\resCostoPequeMax\,ns
con $n\le 10^5$ a hasta \resCostoGrandeMax\,ns con $n=10^6$. El quiebre está
entre $2\cdot 10^5$ y $5\cdot 10^5$ puntos, donde los datos del algoritmo
(unos 60 bytes por punto entre el punto, sus enlaces y su cubeta) dejan de
caber en los 12\,MiB del caché L3, y cada reasignación toca un punto en una
posición aleatoria de la memoria. El \inc{} no sufre ese efecto: \code{std::sort}
y los recorridos acceden a la memoria en orden.

Con \code{int} y todos los puntos ``en la cerradura'', la curva del \ran{} se
dobla hacia abajo desde $n=2\cdot 10^5$: no es el algoritmo sino la entrada, que
pierde vértices al redondear (\cref{sec:enteros}).

\subsection{Porcentaje de puntos en la cerradura (P2 y P3)}

\begin{figure}[!htbp]
\centering
\includegraphics[width=\linewidth]{figuras/porcentaje.pdf}
\caption{Arriba: tiempo contra el porcentaje de puntos en la cerradura
  (circunferencia), mediana de 10 nubes. Abajo: cociente entre las medianas;
  bajo la línea en 1 el \ran{} es más rápido.}
\label{fig:porcentaje}
\end{figure}

La \cref{fig:porcentaje} muestra lo que se esperaba de los mecanismos: el
\inc{} es plano en el porcentaje, y el \ran{} es plano hasta cerca de un 5\,\% y
luego crece, hasta ser \resCirDCociente{} veces más lento con todos los puntos
en la cerradura y $n=10^6$ (\resCirICociente{} con \code{int}).

\paragraph{Puntos aleatorios (P2).} Con $n=10^6$ puntos uniformes la cerradura
tiene una mediana de \resUniH{} vértices (se esperaban unos
\resUniHEsperado), menos de un 0.004\,\%. Los dos algoritmos tardan casi lo
mismo, y cuál gana depende de $n$ (\cref{fig:escalamiento}, abajo a la
izquierda): en tamaños intermedios el \ran{} es algo más rápido (el cociente
$t_\text{ran}/t_\text{inc}$ llega a \resUniCocMinD{} con $n=\resUniCocMinND$ en
\code{double} y a \resUniCocMinI{} con $n=\resUniCocMinNI$ en \code{int}),
pero con $n=10^6$ es más lento: \resUniDRan{} contra \resUniDInc{}\,ms en
\code{double} (cociente \resUniDCociente) y \resUniIRan{} contra
\resUniIInc{}\,ms en \code{int} (\resUniICociente).

\paragraph{El porcentaje de cruce (P3).} La \cref{tab:cruce} interpola, para
cada $n$, el porcentaje en que el cociente cruza 1. El \ran{} es más rápido
solo con $n=10^5$ y menos de un 2--3\,\% de los puntos en la cerradura, y por
poco (un 7\,\% con \code{double}). Con $n=10^4$ y $n=10^6$ el \inc{} es más rápido
en todos los porcentajes. Es la misma ventana de tamaños intermedios que se ve
con los puntos uniformes, medidos en siete tamaños.

\begin{table}[!htbp]
\centering
\caption{Cociente de tiempos $t_\text{ran}/t_\text{inc}$ en los extremos del
  rango de porcentajes, y porcentaje en que cruza 1.}
\label{tab:cruce}
\begin{tabular}{@{}lrrrl@{}}
\toprule
Tipo & $n$ & con 0.1\,\% & con 100\,\% & Cruce \\
\midrule
\tablaCruce
\bottomrule
\end{tabular}
\end{table}

¿Por qué el cruce depende así de $n$? Con un porcentaje pequeño casi todos los
puntos quedan dentro después de pocas inserciones y se descartan: el \ran{} hace
entre \resRporNMin{} y \resRporNMax{} reasignaciones por punto, \emph{sin
importar} $n$, es decir, trabajo lineal. El \inc{} sigue haciendo $n\log n$.
Escribiendo $t_\text{inc}=c_\text{inc}\,n\log_2 n$ y
$t_\text{ran}=c_\text{ran}\,R$, con $R$ las reasignaciones, el cociente se separa
en dos factores:
\[
  \frac{t_\text{ran}}{t_\text{inc}}
  = \underbrace{\frac{R/n}{\log_2 n}}_{\text{operaciones}}\;\cdot\;
    \underbrace{\frac{c_\text{ran}}{c_\text{inc}}}_{\text{costo por operación}}.
\]
La \cref{tab:descomposicion} muestra los dos factores. El de operaciones baja de
$0.32$ a $0.20$ al pasar de $10^4$ a $10^6$, porque $R/n$ es constante y
$\log_2 n$ crece: favorece cada vez más al \ran. Pero una reasignación cuesta
entre 3 y 4 veces lo que cuesta una unidad de $n\log_2 n$ del \inc{} mientras
los datos caben en el caché, y casi 6 veces con $n=10^6$. Con $n=10^5$ el primer
factor ya bajó lo suficiente y el segundo todavía no sube: ahí está la única
ventana en que gana el \ran.

\begin{table}[!htbp]
\centering
\caption{Descomposición del cociente de tiempos con un 0.1\,\% de los puntos
  en la cerradura (medianas de 10 nubes). $c_\text{inc}$ en ns por $n\log_2 n$,
  $c_\text{ran}$ en ns por reasignación.}
\label{tab:descomposicion}
\begin{tabular}{@{}lrrrrrrrr@{}}
\toprule
Tipo & $n$ & $\log_2 n$ & $R/n$ & $\frac{R/n}{\log_2 n}$ & $c_\text{inc}$ &
  $c_\text{ran}$ & $\frac{c_\text{ran}}{c_\text{inc}}$ &
  $\frac{t_\text{ran}}{t_\text{inc}}$ \\
\midrule
\tablaDescomposicion
\bottomrule
\end{tabular}
\end{table}

Entonces el porcentaje de cruce sí se puede determinar, pero no es un número
fijo: depende de $n$ y de la máquina, y aquí existe solo en una ventana
alrededor de $n=10^5$. La descomposición sugiere una hipótesis para otro ciclo
de experimentos (no se midió): con $n$ mucho mayor, cuando todos los accesos
del \ran{} ya fallan en el caché, $c_\text{ran}/c_\text{inc}$ dejaría de crecer
mientras $\log_2 n$ sigue creciendo, y el \ran{} volvería a ganar con
porcentajes pequeños.

\begin{figure}[!htbp]
\centering
\includegraphics[width=\linewidth]{figuras/reasignaciones.pdf}
\caption{Izquierda: reasignaciones por punto del \ran, contra el porcentaje en
  la cerradura; las líneas punteadas son la cota $2\ln n$ de cada $n$.
  Derecha: con coordenadas \code{int}, fracción de los puntos pedidos en la
  circunferencia que siguen siendo vértices después de redondear.}
\label{fig:reasignaciones}
\end{figure}

\subsection{Variabilidad}

\begin{table}[!htbp]
\centering
\caption{Una misma nube de $10^6$ puntos medida con 20 semillas del barajado
  (el \inc{} no usa la semilla: su variación es solo de la máquina). Tiempos en
  ms; CV es el coeficiente de variación.}
\label{tab:varianza}
\small
\setlength{\tabcolsep}{4pt}
\begin{tabular}{@{}lllrrrrr@{}}
\toprule
Tipo & Nube & Algoritmo & Media & Desv. & CV & Reasignaciones & CV \\
\midrule
\tablaVarianza
\bottomrule
\end{tabular}
\end{table}

La \cref{tab:varianza} separa dos fuentes de variación. La del algoritmo se ve
en las reasignaciones, que son exactas: con puntos uniformes varían un 20\,\%
entre semillas (dependen de cuán grande sea la cerradura de los primeros
puntos insertados, que es lo que permite descartar a los demás), y el tiempo
sigue esa variación (correlación \resCorrUniD{} con \code{double},
\resCorrUniI{} con \code{int}). Con todos los puntos en la cerradura las
reasignaciones varían solo un 2\,\%, y la variación del tiempo, de un
10\,\%, no se correlaciona con ellas (\resCorrCirD{} y \resCorrCirI): es de la
máquina. La misma nube con la misma semilla, medida al principio y al final del
experimento, difiere entre \resRuidoMin{} y \resRuidoMax\,\%; el regulador de
frecuencia de un notebook no se puede fijar sin permisos de administrador.
Por eso los resultados se apoyan en medianas de comparaciones pareadas (los dos
algoritmos sobre la misma nube, uno después del otro) y en las cuentas de
operaciones, y una diferencia menor a un 10\,\% entre los algoritmos, como la
ventana de $n=10^5$, debe leerse con cautela.

\subsection{Enteros y reales (P4)}
\label{sec:enteros}

\paragraph{Velocidad.} Con la misma carga de trabajo (puntos uniformes, misma
$h$), \code{int} es más rápido: el \inc{} tarda \resIntDoubleUniInc{} y el \ran{}
\resIntDoubleUniRan{} veces lo que tardan con \code{double} con $n=10^6$. Una
razón es la memoria: \code{Punto<T>} guarda también la coordenada $z$ de la
Tarea~1, así que un \code{Punto<int>} ocupa 12 bytes y un \code{Punto<double>}
24, y ambos algoritmos mueven puntos completos.

\paragraph{Exactitud.} Con \code{int} todas las orientaciones son exactas (por
el rango de coordenadas de los supuestos), y las dos cerraduras coincidieron en
todos los casos, también con puntos colineales. Con \code{double} solo
coinciden mientras ninguna orientación dudosa decida algo
(\cref{sec:colineal}).

\paragraph{El redondeo cambia la entrada.} Al redondear puntos de la
circunferencia, dos vecinos quedan separados de la cuerda de sus vecinos por
apenas $R(1-\cos\frac{2\pi}{k})\approx 2\pi^2R/k^2$, y el redondeo los mueve
hasta medio entero. Con $R=10^8$ esa separación baja de 1 cuando
$k>\pi\sqrt{2R}\approx 44\,000$, y entonces muchos puntos dejan de estar en
posición convexa. La \cref{fig:reasignaciones} (derecha) lo confirma: las
pérdidas empiezan con $k=50\,000$ (5\,\% de $10^6$ y 50\,\% de $10^5$), y con
$n=10^6$ y 100\,\% pedido solo quedan \resIntTodosH{} vértices. No es un
problema del generador sino de los enteros: en una grilla de $N\times N$ caben
a lo más del orden de $N^{2/3}$ puntos en posición estrictamente convexa
\cite{jarnik}. Con \code{int} el porcentaje real se mide en la salida y es el
que corresponde usar.

\subsection{Puntos colineales (P5)}
\label{sec:colineal}

\begin{table}[!htbp]
\centering
\caption{Cerraduras con un 10\,\% de los puntos en el borde, 10 nubes por
  fila: porcentaje de nubes con cerraduras idénticas y rango del número de
  vértices de cada algoritmo.}
\label{tab:colineal}
\begin{tabular}{@{}llrrrr@{}}
\toprule
Nube & Tipo & $n$ & Idénticas & $h$ \inc & $h$ \ran \\
\midrule
\tablaColinealidad
\bottomrule
\end{tabular}
\end{table}

La \cref{tab:colineal} tiene tres situaciones:

\begin{itemize}
  \item \textbf{Colinealidad exacta} (cuadrado sin rotar, \code{int} y
        \code{double}): los puntos de los lados tienen exactamente la misma $x$ o
        la misma $y$, las orientaciones dan 0 sin error, y ambos algoritmos
        entregan las 4 esquinas. Las reglas de la \cref{sec:supuestos} (sacar con
        orientación $\ge 0$; prueba estricta y recorrido no estricto) bastan.
  \item \textbf{Sin colineales} (circunferencia): siempre coinciden, también con
        \code{double}, porque los vértices están lejos de ser colineales
        comparados con el error de redondeo.
  \item \textbf{Colinealidad que el tipo no representa} (cuadrado rotado en
        $30^\circ$): con \code{int} los puntos redondeados ya no son
        colineales, pero las orientaciones son exactas y ambos algoritmos
        calculan la misma cerradura, la de los puntos enteros que sobresalen de
        cada lado (25 a 58 vértices). Con \code{double} los puntos quedan a un
        redondeo de la recta, las orientaciones entre ellos son del tamaño del
        error, y las cerraduras \emph{nunca} coinciden.
\end{itemize}

Para saber cuál de las dos es la correcta en el último caso, se recalculó la
cerradura de los mismos \code{double} con aritmética racional exacta
(\code{analisis/exactitud.py}, con \code{fractions.Fraction}; cada
\code{double} es un racional). La \cref{tab:exactitud} muestra que
\emph{ninguno} de los dos algoritmos entrega la cerradura exacta, y que fallan de
maneras opuestas. El \inc{} agrega vértices que no lo son: cuando el signo de
una orientación casi nula sale mal, cree ver una vuelta a la derecha y deja en
la cadena un punto que debía sacar. El \ran{} pierde vértices verdaderos: el
mismo error en la prueba ``está fuera'' descarta un punto, y el descarte es
para siempre. Además, como el resultado del \ran{} depende del orden de
inserción, con \code{double} cambia con la semilla.

\begin{table}[!htbp]
\centering
\caption{Cuadrado rotado en $30^\circ$ con \code{double}, comparado con la
  cerradura exacta de los mismos puntos (10 nubes por fila): nubes en que la
  cerradura calculada es la exacta, y promedio de vértices exactos que faltan y
  de vértices calculados que sobran.}
\label{tab:exactitud}
\begin{tabular}{@{}rrrrrrrr@{}}
\toprule
& & \multicolumn{3}{c}{\inc} & \multicolumn{3}{c}{\ran} \\
\cmidrule(lr){3-5}\cmidrule(l){6-8}
$n$ & $h$ exacta & Exactas & Faltan & Sobran & Exactas & Faltan & Sobran \\
\midrule
\tablaExactitud
\bottomrule
\end{tabular}
\end{table}

Como pide el enunciado, los algoritmos no se hicieron consistentes en este caso:
hacerlo requiere predicados exactos (por ejemplo, los adaptativos de Shewchuk) o
coordenadas enteras. Con los supuestos de este trabajo, la respuesta práctica
es usar coordenadas enteras cuando la entrada tiene colinealidades.

\section{Conclusiones}

\begin{enumerate}[label=\textbf{P\arabic*}]
  \item Ambos algoritmos hacen $\Theta(n\log n)$ operaciones (el \ran{} en su peor
        caso, con todos los puntos en la cerradura). El \inc{} también lo
        cumple en tiempo, con unos 7\,ns por $n\log_2 n$; el \ran{} se aleja de
        $n\log n$ en tiempo porque cada operación se encarece cuando los datos
        no caben en el caché.
  \item Con puntos uniformes los dos tienen tiempos parecidos: el \ran{} es
        hasta un 10--16\,\% más rápido en tamaños intermedios y un
        14--16\,\% más lento con $n=10^6$.
  \item El \inc{} no depende del porcentaje de puntos en la cerradura; el \ran{}
        hace trabajo lineal hasta un 5\,\% y luego crece hasta ser casi 10 veces
        más lento con el 100\,\%. El \ran{} gana solo en tamaños intermedios
        (aquí, $n=10^5$) y con menos de un 2--3\,\% en la cerradura; el
        porcentaje de cruce se puede determinar
        experimentalmente, pero depende de $n$ y de la máquina, por la
        competencia entre $\log n$ y el costo de memoria
        (\cref{tab:descomposicion}).
  \item Con la misma entrada, \code{int} es un 10--12\,\% más rápido y hace
        exactas todas las decisiones. Pero el redondeo cambia la entrada: al
        redondear una circunferencia de radio $10^8$, desde unos $44\,000$
        puntos estos dejan de estar en posición convexa.
  \item Con colinealidad exacta o sin colineales, ambas cerraduras coinciden
        siempre; con puntos casi colineales en \code{double} nunca coinciden,
        y ninguna es la exacta: el \inc{} agrega vértices y el \ran{} los
        pierde.
\end{enumerate}

En la práctica, el \inc{} es la mejor opción general: es más simple, su tiempo
no depende de la entrada y es más rápido en casi todos los casos medidos.



\begin{thebibliography}{9}
\bibitem{cmu} G.~Miller, J.~Li. \emph{Lecture Notes: Quicksort and Random
  Incremental Convex Hull}. 15-451 Algorithms, Carnegie Mellon University.
  \url{https://www.cs.cmu.edu/~15451-f20/LectureNotes/2D-ConvexHull.pdf}
\bibitem{deberg} M.~de Berg, O.~Cheong, M.~van Kreveld, M.~Overmars.
  \emph{Computational Geometry: Algorithms and Applications}, 3.\textsuperscript{a}
  ed. Springer, 2008. Sección~1.1.
\bibitem{moret} B.~M.~E.~Moret. \emph{Towards a Discipline of Experimental
  Algorithmics}. DIMACS Series in Discrete Mathematics and Theoretical
  Computer Science, vol.~59, 2002.
\bibitem{renyi} A.~Rényi, R.~Sulanke. Über die konvexe Hülle von $n$ zufällig
  gewählten Punkten. \emph{Z. Wahrscheinlichkeitstheorie}, 2:75--84, 1963.
\bibitem{jarnik} V.~Jarník. Über die Gitterpunkte auf konvexen Kurven.
  \emph{Math. Z.}, 24:500--518, 1926.
\end{thebibliography}

\end{document}
